\chapter{DeFi Credit and Leverage}\label{ch:m3:defi-credit-and-leverage}

On 11 October 2022 a trader bought a thin \omterm{def:m3:blockchains-for-traders:key}{token} on the few exchanges whose prices fed a lending venue's \omterm{def:m3:blockchains-for-traders:contract}{oracle}.
Within half an hour the price the \omterm{def:m3:blockchains-for-traders:contract}{oracle} reported rose more than thirteenfold. The trader's position in the
venue, now marked at that price, was worth enough to serve as collateral for borrowing everything the venue
held: he borrowed and withdrew more than USD~110 million of other \omterm{def:m3:blockchains-for-traders:key}{tokens} and left the inflated collateral
behind. Lending on a \omterm{def:m3:blockchains-for-traders:chain}{blockchain} works without a credit officer: a program lends against collateral at a
loan-to-value it knows, liquidates anyone whose collateral falls below a threshold, and prices everything
with an \omterm{def:m3:blockchains-for-traders:contract}{oracle}. This chapter covers the pools that do the lending, the liquidations that keep them solvent, the
\omterm{def:m3:defi-credit-and-leverage:flash}{flash loans} that exist nowhere else, the \omterm{def:m3:defi-credit-and-leverage:staking}{staking yield} on which much of the collateral earns a return, and
the two ways the system has broken: \omterm{def:m3:blockchains-for-traders:stable}{stablecoins} whose peg was an algorithm, and \omterm{def:m3:blockchains-for-traders:contract}{oracles} that could be
bought.

\section{Lending pools}

\begin{definition}[Lending pool, kinked interest-rate curve]\label{def:m3:defi-credit-and-leverage:pool}
A \emph{lending pool}\index{lending pool} is a \omterm{def:m3:blockchains-for-traders:contract}{smart contract} that takes deposits of a \omterm{def:m3:blockchains-for-traders:key}{token} from suppliers and
lends them to borrowers who post other \omterm{def:m3:blockchains-for-traders:key}{tokens} as collateral, paying suppliers the interest borrowers pay less a
share kept by the protocol. A \emph{kinked interest-rate curve}\index{kinked interest-rate curve} sets the borrow
rate as a function of the pool's utilisation, rising gently up to a target utilisation and steeply beyond it.
\end{definition}

Utilisation is the share of supplied \omterm{def:m3:blockchains-for-traders:key}{tokens} that is lent out, as in the securities lending of One Quant Book 1,
chapter 16. The pool has no maturity: suppliers can withdraw at any time, but only from what is not lent. The
kink is the pool's defence of that liquidity. Above it, borrowing becomes so expensive that borrowers repay and
suppliers arrive, pulling utilisation back and leaving \omterm{def:m3:blockchains-for-traders:key}{tokens} for those who want to withdraw.

\begin{proposition}[Supply rate]\label{prop:m3:defi-credit-and-leverage:supply}
With utilisation $u$, borrow rate $r_b(u)$ and a reserve factor $f_{\mathrm{res}}$ kept by the protocol, suppliers earn
$r_s = r_b(u)\,u\,(1 - f_{\mathrm{res}})$.
\end{proposition}

\begin{proof}
Borrowers pay $r_b$ on the lent amount, a fraction $u$ of the supply; the protocol keeps $f_{\mathrm{res}}$ of it.
\end{proof}

\begin{omfigure}
\begin{tikzpicture}
\begin{axis}[width=11cm,height=5cm,
  xlabel={utilisation, \%},ylabel={rate, \% a year},
  tick label style={font=\small},label style={font=\small},
  xmin=0,xmax=100,ymin=0,ymax=65,grid=major,grid style={gray!25},
  legend columns=2,legend style={font=\footnotesize,at={(0.5,-0.3)},anchor=north,draw=none,fill=none,/tikz/every even column/.append style={column sep=8pt}},
  table/col sep=comma]
\addplot[omThm,very thick] table[x=u,y=borrow]{figdata/markets-3/23-defi-credit-and-leverage/rates.csv};\addlegendentry{borrow rate}
\addplot[omDef,very thick,dashed] table[x=u,y=supply]{figdata/markets-3/23-defi-credit-and-leverage/rates.csv};\addlegendentry{supply rate}
\end{axis}
\end{tikzpicture}

\omcaption{A \omterm{def:m3:defi-credit-and-leverage:pool}{kinked interest-rate curve}: 4\% at the 90\% optimal utilisation, rising by 60 points more to full utilisation, and the
supply rate it implies with a 10\% reserve factor (\cref{prop:m3:defi-credit-and-leverage:supply}). Illustrative parameters of the
two-slope form. Data: the chapter's tutorial.}\label{fig:m3:defi-credit-and-leverage:rates}
\end{omfigure}

\begin{dated}{2026-09}{A lending protocol's rate model and flash-loan fee}\label{dat:m3:defi-credit-and-leverage:aave}
Aave's documentation describes its v3 interest-rate strategy as based on two slopes, one below an optimal usage ratio and another from
that point to 100\%, with a base variable borrow rate, per reserve. Positions are over-collateralised, risk is tracked by a \omterm{def:m3:defi-credit-and-leverage:ltv}{health factor}
with per-reserve liquidation thresholds, and a position with a \omterm{def:m3:defi-credit-and-leverage:ltv}{health factor} below 1 can be liquidated by anyone, who repays part of the
debt and receives collateral at a discount, the \omterm{def:m3:defi-credit-and-leverage:ltv}{liquidation bonus}. In the v3 contracts a liquidation may repay at most 50\% of the debt
(the close factor), and all of it when the \omterm{def:m3:defi-credit-and-leverage:ltv}{health factor} is at or below 0.95 or the position is below USD~2\,000. The flash-loan fee was set at 0.05\%
at deployment and can be changed by governance.
\end{dated}

\section{Liquidations}

\begin{definition}[Loan-to-value ratio, health factor, liquidation bonus]\label{def:m3:defi-credit-and-leverage:ltv}
The \emph{loan-to-value ratio}\index{loan-to-value ratio} of a collateral asset is the maximum a borrower may borrow per dollar of it. The
\emph{health factor}\index{health factor} of an account is the sum over its collateral of value times liquidation threshold, divided by the
value of its debt; below one, the account may be liquidated. The \emph{liquidation bonus}\index{liquidation bonus} is the discount at which a
liquidator receives collateral for the debt it repays.
\end{definition}

The \omterm{def:m3:defi-credit-and-leverage:ltv}{loan-to-value ratio} is a haircut (One Quant Book 1, chapter 6) applied when borrowing; the liquidation threshold, a little higher, is the
haircut at which the position is taken. There is no margin call: the borrower is not asked for more collateral, it is liquidated by whoever
acts first, and the bonus pays them to act. A borrower with 100 ether of collateral at a threshold of 82.5\% and USD~240\,000 of debt has a
\omterm{def:m3:defi-credit-and-leverage:ltv}{health factor} of 1.03 at 3\,000 dollars and is liquidated below 2\,909.

\begin{proposition}[Liquidation arithmetic]\label{prop:m3:defi-credit-and-leverage:liq}
A liquidator repays an amount $R$ of debt, at most the close factor $c$ times the debt, and receives collateral worth $R(1 + \beta)$ at the
\omterm{def:m3:blockchains-for-traders:contract}{oracle} price, with $\beta$ the bonus. The borrower's \omterm{def:m3:defi-credit-and-leverage:ltv}{health factor} after a liquidation of $R$ is
$(C_\theta - R(1 + \beta)\theta)/(D - R)$, where $C_\theta$ is the threshold-weighted collateral value, $D$ the debt value and $\theta$ the
collateral's threshold; it rises only if $(1 + \beta)\theta < C_\theta/D$, which holds when the \omterm{def:m3:defi-credit-and-leverage:ltv}{health factor} was above $(1 + \beta)\theta$.
\end{proposition}

\begin{proof}
The debt falls by $R$ and the threshold-weighted collateral by $R(1 + \beta)\theta$. The ratio $(C_\theta - xR)/(D - R)$ with $x = (1 +
\beta)\theta$ exceeds $C_\theta/D$ exactly when $C_\theta/D > x$.
\end{proof}

\begin{omfigure}
\begin{tikzpicture}
\begin{axis}[width=11cm,height=5cm,x dir=reverse,
  xlabel={ether price, USD},ylabel={health factor},
  tick label style={font=\small},label style={font=\small},
  xmin=2000,xmax=3300,ymin=0.6,ymax=1.2,grid=major,grid style={gray!25},
  xtick={2000,2250,2500,2750,3000,3250},xticklabels={2\,000,2\,250,2\,500,2\,750,3\,000,3\,250},scaled x ticks=false,
  table/col sep=comma]
\addplot[omThm,very thick] table[x=price,y=hf]{figdata/markets-3/23-defi-credit-and-leverage/health.csv};
\draw[omDef,thick,dashed] (axis cs:3300,1) -- (axis cs:2000,1);
\draw[omExo,thick] (axis cs:2909,0.6) -- (axis cs:2909,1.2);
\node[font=\footnotesize,anchor=west,text=omExo] at (axis cs:2890,0.68) {liquidation below 2\,909};
\end{axis}
\end{tikzpicture}

\omcaption{\omterm{def:m3:defi-credit-and-leverage:ltv}{Health factor} of a borrower with 100 ether of collateral (threshold 82.5\%) and USD~240\,000 of debt as the price of ether falls.
Data: the chapter's tutorial.}\label{fig:m3:defi-credit-and-leverage:health}
\end{omfigure}

The proposition has a sting. A position that falls far below one before anyone acts, because the price gapped or the chain was congested, can
be pushed deeper by its own liquidation: each dollar repaid removes more than a dollar of threshold-weighted collateral. When the collateral left
is worth less than the debt, the pool is left with bad debt that its suppliers bear.

\section{Flash loans}

\begin{definition}[Flash loan]\label{def:m3:defi-credit-and-leverage:flash}
A \emph{flash loan}\index{flash loan} is a loan without collateral that must be repaid, with a fee, within the same transaction; if it is not,
the whole transaction reverts and the loan never happened.
\end{definition}

A \omterm{def:m3:defi-credit-and-leverage:flash}{flash loan} has no credit risk for the lender because the chain's atomicity enforces repayment: the transaction either ends with the pool repaid
or does not exist. It lets anyone act with a pool's capital for one transaction. The canonical use is a liquidation by a liquidator with no capital:
borrow the \omterm{def:m3:blockchains-for-traders:stable}{stablecoins} to repay, receive the collateral with its bonus, sell it on an \omterm{def:m3:automated-market-makers:amm}{automated market maker}, repay the \omterm{def:m3:defi-credit-and-leverage:flash}{flash loan}, keep the rest.
In the tutorial, liquidating half of the borrower's debt at an ether price of 2\,800 seizes 45 ether, sells them for USD~124\,505 in a pool of
5\,000 ether, repays the USD~120\,000 borrowed and USD~60 of fee, and leaves USD~4\,445. The same atomicity serves attacks: a \omterm{def:m3:defi-credit-and-leverage:flash}{flash loan} can
move a price, use the moved price, and restore it, all in one transaction. Qin, Zhou, Livshits and Gervais analysed two such attacks, both
in February 2020, and measured returns on investment above 500\,000\%.

\begin{omfigure}
\centering
\begin{tikzpicture}[font=\small,
  s/.style={draw,thick,rounded corners=2pt,align=center,minimum height=0.9cm,text width=2.2cm},
  arr/.style={-{Stealth[length=2.2mm]},thick}]
\node[s,draw=omDef,fill=omDef!8] (a) at (0,0) {flash-borrow\\120\,000 USDC};
\node[s,draw=omThm,fill=omThm!8] (b) at (2.8,0) {repay half\\the debt};
\node[s,draw=omThm,fill=omThm!8] (c) at (5.6,0) {receive 45 ETH\\(5\% bonus)};
\node[s,draw=omProp,fill=omProp!8] (d) at (8.4,0) {sell in a pool:\\124\,505 USDC};
\node[s,draw=omDef,fill=omDef!8] (e) at (11.2,0) {repay 120\,060;\\keep 4\,445};
\draw[arr] (a) -- (b); \draw[arr] (b) -- (c); \draw[arr] (c) -- (d); \draw[arr] (d) -- (e);
\draw[thick,gray,dashed,rounded corners] (-1.3,-0.8) rectangle (12.5,0.8);
\node[font=\footnotesize,gray,anchor=north] at (5.6,-0.85) {one transaction: if the last step fails, none of it happened};
\end{tikzpicture}

\omcaption{A liquidation funded by a \omterm{def:m3:defi-credit-and-leverage:flash}{flash loan}, with the tutorial's numbers. Schematic.}\label{fig:m3:defi-credit-and-leverage:flash}
\end{omfigure}

\section{Liquid staking, restaking and the staking yield}

\begin{definition}[Staking yield, liquid staking token, restaking]\label{def:m3:defi-credit-and-leverage:staking}
The \emph{staking yield}\index{staking yield} is the return earned by locking a proof-of-stake chain's \omterm{def:m3:blockchains-for-traders:key}{token} with a \omterm{def:m3:blockchains-for-traders:pow}{validator}
(\cref{ch:m3:blockchains-for-traders}): new issuance and a share of fees, less penalties. A \emph{liquid staking token}\index{liquid staking token} is a
transferable \omterm{def:m3:blockchains-for-traders:key}{token} that represents staked \omterm{def:m3:blockchains-for-traders:key}{tokens} and their accruing rewards, issued by a protocol that stakes deposits on its users' behalf.
\emph{Restaking}\index{restaking} is committing already-staked \omterm{def:m3:blockchains-for-traders:key}{tokens}, or liquid staking \omterm{def:m3:blockchains-for-traders:key}{tokens}, as security for further services, for additional
rewards and additional risk of their being destroyed for misbehaviour.
\end{definition}

\begin{dated}{2026-09}{A staking yield}\label{dat:m3:defi-credit-and-leverage:lido}
Lido's public API reported a seven-day moving average annual percentage rate of 2.24\% for its staked-ether \omterm{def:m3:blockchains-for-traders:key}{token}, stETH, on 24 September 2026,
and 9.76 million ether staked through it: about 22\% of the 43.8 million ether held by \omterm{def:m3:blockchains-for-traders:pow}{validators} on the beacon chain that day.
\end{dated}

The \omterm{def:m3:defi-credit-and-leverage:staking}{staking yield} is crypto's nearest thing to a risk-free rate for ether, and the \omterm{def:m3:defi-credit-and-leverage:staking}{liquid staking token} is the collateral on which much of the
leverage is built: deposit stETH, borrow ether, stake it, repeat, earning the \omterm{def:m3:defi-credit-and-leverage:staking}{staking yield} on several times the capital as long as the borrow rate
stays below it. The loop's risks are the ones of any carry trade funded short: the borrow rate can rise above the yield when utilisation climbs,
and the \omterm{def:m3:defi-credit-and-leverage:staking}{liquid staking token} can trade below the value of the ether behind it when holders want to leave faster than withdrawals from staking allow,
triggering liquidations of the loopers. \omterm{def:m3:defi-credit-and-leverage:staking}{Restaking} layers further claims on the same collateral: EigenLayer's contracts accept \omterm{def:m3:defi-credit-and-leverage:staking}{liquid staking
tokens}, beacon-chain ether and other \omterm{def:m3:blockchains-for-traders:key}{tokens}, and let each operator allocate a share of its delegated stake to be slashable by a given service.

\section{Algorithmic stablecoins and oracle manipulation}

\begin{definition}[Algorithmic stablecoin]\label{def:m3:defi-credit-and-leverage:algo}
An \emph{algorithmic stablecoin}\index{algorithmic stablecoin} is a \omterm{def:m3:blockchains-for-traders:key}{token} that aims at a fixed value without holding reserves of that value, relying
instead on a mechanism that exchanges it for another \omterm{def:m3:blockchains-for-traders:key}{token} of floating value created or destroyed to absorb demand.
\end{definition}

The mechanism holds the peg only while the floating \omterm{def:m3:blockchains-for-traders:key}{token} is worth more than the \omterm{def:m3:blockchains-for-traders:stable}{stablecoins} outstanding; when confidence falls, redemptions create
floating \omterm{def:m3:blockchains-for-traders:key}{tokens} faster than anyone wants them, their price falls, and each redemption creates more of them. According to the SEC, Terraform Labs
marketed its \omterm{def:m3:blockchains-for-traders:stable}{stablecoin} UST, which was to hold its dollar peg by being exchangeable for its \omterm{def:m3:blockchains-for-traders:key}{token} LUNA, as ``yield-bearing'', paying as much as 20\%
through its Anchor protocol; in May 2022 UST lost its peg and it and its sister \omterm{def:m3:blockchains-for-traders:key}{tokens} fell close to zero.

\begin{definition}[Oracle manipulation]\label{def:m3:defi-credit-and-leverage:oracle}
\emph{Oracle manipulation}\index{oracle manipulation} is moving the price that an \omterm{def:m3:blockchains-for-traders:contract}{oracle} reports to a \omterm{def:m3:blockchains-for-traders:contract}{smart contract}, usually by trading in the thin
markets it reads, in order to borrow against inflated collateral, avoid a liquidation or trigger one.
\end{definition}

\begin{dated}{2026-09}{The Mango Markets case}\label{dat:m3:defi-credit-and-leverage:mango}
The CFTC alleged in January 2023 that on 11 October 2022 Avraham Eisenberg, through two accounts on Mango Markets, built large leveraged positions in a
swap on the MNGO \omterm{def:m3:blockchains-for-traders:key}{token}, pumped MNGO's price on three exchanges feeding the venue's \omterm{def:m3:blockchains-for-traders:contract}{oracle}, which reported a more than thirteenfold rise in 30 minutes,
and used the inflated positions as collateral to withdraw over USD~110 million, of which about USD~67 million was later returned. A jury convicted him of
commodities fraud, commodities manipulation and wire fraud; on 23 May 2025 the trial judge granted his motion under Rule 29, vacating the first two
counts and entering a judgment of acquittal on the third.
\end{dated}

\begin{proposition}[The pump that borrows a pool dry]\label{prop:m3:defi-credit-and-leverage:pump}
A pool lends up to $L$ against collateral of $q$ \omterm{def:m3:blockchains-for-traders:key}{tokens} at price $P$ and loan-to-value $\lambda$, with the price read from a constant-product market
holding $y$ dollars. Pushing that market's price up by a factor $k$ costs $y(\sqrt{k} - 1)$ dollars and buys a fraction $1 - 1/\sqrt{k}$ of its \omterm{def:m3:blockchains-for-traders:key}{tokens}.
The attacker can borrow the whole pool once $k \ge k^* = L/(\lambda qP)$; its gain from borrowing and abandoning the collateral is
$\min(L, \lambda qPk) - qP - y(\sqrt{k} - 1)$ plus whatever the \omterm{def:m3:blockchains-for-traders:key}{tokens} bought are worth afterwards.
\end{proposition}

\begin{proof}
Keeping $xy$ constant while the price $y/x$ rises by $k$ multiplies $y$ by $\sqrt k$ and divides $x$ by $\sqrt k$. Borrowing power is $\lambda qPk$, capped
by the pool; the collateral is left behind.
\end{proof}

\begin{omfigure}
\begin{tikzpicture}
\begin{axis}[width=11cm,height=5cm,
  xlabel={pump of the oracle's source price, multiple},ylabel={attacker's gain, USD million},
  tick label style={font=\small},label style={font=\small},
  xmin=1,xmax=40,ymin=-5,ymax=45,grid=major,grid style={gray!25},
  table/col sep=comma]
\addplot[omThm,very thick] table[x=k,y=gain]{figdata/markets-3/23-defi-credit-and-leverage/pump.csv};
\draw[omDef,thick,dashed] (axis cs:16.67,-5) -- (axis cs:16.67,45);
\node[font=\footnotesize,anchor=west,text=omDef] at (axis cs:17.2,5) {$k^* = 16.7$: pool borrowed dry};
\end{axis}
\end{tikzpicture}

\omcaption{The \omterm{def:m3:blockchains-for-traders:contract}{oracle} pump of \cref{prop:m3:defi-credit-and-leverage:pump}: 5 million \omterm{def:m3:blockchains-for-traders:key}{tokens} at USD~1 as collateral (loan-to-value 60\%), a pool of USD~50
million, and a source market holding USD~2 million; \omterm{def:m3:blockchains-for-traders:key}{tokens} bought in the pump are assumed worthless afterwards. The gain turns positive at a pump of
less than two times, peaks when the pool is empty, and then falls as further pumping only costs. Illustrative parameters. Data: the chapter's tutorial.}\label{fig:m3:defi-credit-and-leverage:pump}
\end{omfigure}

The defences follow from the formula: \omterm{def:m3:defi-credit-and-leverage:ltv}{loan-to-value ratios} that fall with the thinness of the collateral's markets, caps on how much can be borrowed
against any one collateral, \omterm{def:m3:blockchains-for-traders:contract}{oracles} that read deep markets and smooth over time, and circuit breakers on sudden price moves. None of them is free:
each makes the pool less useful to honest borrowers.

\section{Tutorial: a lending pool under stress}\label{tut:m3:defi-credit-and-leverage:pool}

\begin{tutorial}
\textbf{Goal.} Simulate a \omterm{def:m3:defi-credit-and-leverage:pool}{lending pool}: the borrow rate as a kinked function of utilisation, a borrower's \omterm{def:m3:defi-credit-and-leverage:ltv}{health factor} through a price fall, a
liquidation with close factor and bonus, and a liquidation funded by a \omterm{def:m3:defi-credit-and-leverage:flash}{flash loan}. \textbf{End state:} the chapter's figures and the numbers of the
text.
\begin{tutsteps}
\item \textbf{Liquidation and the \omterm{def:m3:defi-credit-and-leverage:flash}{flash loan}.} The engine's two most delicate functions.
\omcode{code/firm/lendpool/firm_lendpool.py}{93}{119}{Liquidation with close factor and bonus, and a flash loan that reverts unless repaid.}\label{lst:m3:defi-credit-and-leverage:liq}
\item \textbf{The flash liquidation.} Borrow, liquidate, sell the collateral in a pool with \texttt{firm.amm}, repay.
\omcode{code/markets-3/23-defi-credit-and-leverage/python/m3_defi.py}{38}{57}{A liquidation funded by a flash loan.}\label{lst:m3:defi-credit-and-leverage:flash}
\item \textbf{Run} \texttt{flash\_liquidation()}, \texttt{pump\_attack} over pumps from 1 to 40, and \texttt{fig\_defi.py}.
\end{tutsteps}
\textbf{What to change next.} Let the price gap to 2\,400 before the liquidation and find the pool's bad debt; add a second liquidator competing in the same
block; make the \omterm{def:m3:blockchains-for-traders:contract}{oracle} a 30-minute average and recompute the pump's cost.
\end{tutorial}

\section{Build: the lending pool engine}\label{bld:m3:defi-credit-and-leverage:lendpool}

\begin{build}
\bfield{Purpose} The miniature firm lends, borrows and liquidates on chain: it must know its \omterm{def:m3:defi-credit-and-leverage:ltv}{health factors}, the rates it will pay as utilisation moves,
the profit of a liquidation after the sale of the collateral, and the exposure of a pool it supplies to bad debt.
\bfield{Interface} \texttt{RateModel(base, slope1, slope2, optimal)}; \texttt{Reserve} with \texttt{utilisation}, \texttt{rates}, \texttt{accrue},
\texttt{available}; \texttt{Pool} with \texttt{supply}, \texttt{borrow}, \texttt{health}, \texttt{liquidate}, \texttt{flash\_loan}. Collateral sales through
\texttt{firm.amm}.
\bfield{Rules} Debt and deposits held as scaled balances against indices; borrowing refused above the loan-to-value or the available liquidity; liquidation
only below a \omterm{def:m3:defi-credit-and-leverage:ltv}{health factor} of one and at most the close factor; a \omterm{def:m3:defi-credit-and-leverage:flash}{flash loan} that is not repaid raises and leaves no trace.
\bfield{Acceptance tests} \texttt{code/firm/lendpool/tests/}: the kink; \omterm{def:m3:defi-credit-and-leverage:ltv}{health factor} and accrual; refusal above the loan-to-value; liquidation of a healthy
account refused; seizure with the bonus; a \omterm{def:m3:defi-credit-and-leverage:flash}{flash loan}'s fee and its revert.
\bfield{Stretch} Several collateral assets with different thresholds; bad-debt accounting; isolation modes and borrow caps; the rate's reaction to a run on
deposits.
\end{build}

\begin{omsources}
\item Aave documentation: v3 overview, interest-rate strategy, flash loans; the contract LiquidationLogic in the aave-v3-origin repository. Accessed September 2026.
\item EigenLayer core contracts documentation (Layr-Labs/eigenlayer-contracts, docs/README), September 2026.
\item K.~Qin, L.~Zhou, B.~Livshits and A.~Gervais, ``Attacking the DeFi Ecosystem with Flash Loans for Fun and Profit'', Financial Cryptography 2021
(arXiv:2003.03810).
\item CFTC, press release 8647-23 (charges against Avraham Eisenberg), January 2023; United States v.\ Eisenberg (S.D.N.Y.), opinion and order of 23
May 2025 (Doc 220).
\item SEC, press release 2023-32, ``SEC Charges Terraform and CEO Do Kwon with Defrauding Investors'', 16 February 2023.
\item Lido, public stETH APR and statistics endpoints; beacon-chain validator balances from a public beacon node; 24 September 2026.
\end{omsources}

\section{Exercises}

\begin{exercise}[$\star$]\label{exo:m3:defi-credit-and-leverage:1}
With the chapter's rate curve, what are the borrow and supply rates at 50\% and at 95\% utilisation?
\end{exercise}

\begin{exercise}[$\star$]\label{exo:m3:defi-credit-and-leverage:2}
A borrower has 50 ether of collateral (threshold 82.5\%) and USD~100\,000 of debt. At what ether price is it liquidated?
\end{exercise}

\begin{exercise}[$\star$]\label{exo:m3:defi-credit-and-leverage:3}
Why does a \omterm{def:m3:defi-credit-and-leverage:flash}{flash loan} carry no credit risk for the pool?
\end{exercise}

\begin{exercise}[$\star\star$]\label{exo:m3:defi-credit-and-leverage:4}
A liquidator repays USD~50\,000 of debt against ether at 2\,500 with a 5\% bonus. How much ether does it receive?
\end{exercise}

\begin{exercise}[$\star\star$]\label{exo:m3:defi-credit-and-leverage:5}
A borrower's \omterm{def:m3:defi-credit-and-leverage:ltv}{health factor} is 0.80 and the collateral's threshold is 82.5\% with a 10\% bonus. Does liquidation improve its \omterm{def:m3:defi-credit-and-leverage:ltv}{health factor}?
\end{exercise}

\begin{exercise}[$\star\star$]\label{exo:m3:defi-credit-and-leverage:6}
Explain the looping trade on a \omterm{def:m3:defi-credit-and-leverage:staking}{liquid staking token} and name two ways it loses money.
\end{exercise}

\begin{exercise}[$\star\star\star$]\label{exo:m3:defi-credit-and-leverage:7}
\emph{Coding.} With \texttt{pump\_attack}, find the smallest pump at which the attack pays, with and without the \omterm{def:m3:blockchains-for-traders:key}{tokens} bought keeping a value of USD~1.
\end{exercise}

\begin{exercise}[$\star\star\star$]\label{exo:m3:defi-credit-and-leverage:8}
\emph{Find the flaw.} ``Our pool is safe: every loan is over-collateralised.''
\end{exercise}

\section{Problem: The Oracle Pump}

\begin{problem}[{Weekend problem --- borrowing a pool dry}]\label{pb:m3:defi-credit-and-leverage:1}
A \omterm{def:m3:defi-credit-and-leverage:pool}{lending pool} holds USD~50 million of \omterm{def:m3:blockchains-for-traders:stable}{stablecoins} and accepts a thin \omterm{def:m3:blockchains-for-traders:key}{token} as collateral at a loan-to-value of 60\%, priced by an \omterm{def:m3:blockchains-for-traders:contract}{oracle} that reads a
constant-product market holding USD~2 million and 2 million \omterm{def:m3:blockchains-for-traders:key}{tokens} (price USD~1). An attacker holds 5 million \omterm{def:m3:blockchains-for-traders:key}{tokens}.

\medskip\noindent\textbf{Part I --- The mechanics.}
\begin{enumerate}
\item What can the attacker borrow honestly?
\item By what factor must the \omterm{def:m3:blockchains-for-traders:contract}{oracle}'s price rise for the attacker to borrow the whole pool?
\item What does pushing the market's price up by that factor cost, and how many \omterm{def:m3:blockchains-for-traders:key}{tokens} does it buy?
\item What is the attacker's gain if the \omterm{def:m3:blockchains-for-traders:key}{tokens} bought become worthless? If they keep USD~1?
\item Why does the gain fall for pumps beyond that factor?
\end{enumerate}

\medskip\noindent\textbf{Part II --- The thresholds.}
\begin{enumerate}[resume]
\item What is the smallest pump that pays?
\item How does the answer change if the source market holds USD~20 million?
\item What if the loan-to-value is 30\%?
\item What borrow cap on this collateral would limit the attacker's gain to zero at any pump?
\item What would a 30-minute time-weighted \omterm{def:m3:blockchains-for-traders:contract}{oracle} change?
\end{enumerate}

\medskip\noindent\textbf{Part III --- The case.}
\begin{enumerate}[resume]
\item What did the CFTC allege happened on Mango Markets?
\item What did the trial judge decide after the jury's verdict, and on what kind of motion?
\item Who bears the loss when a pool is borrowed dry?
\item How does this attack differ from a flash-loan attack?
\item What would a lending protocol's governance do next?
\end{enumerate}

\medskip\noindent\textbf{Part IV --- Judgement.}
\begin{enumerate}[resume]
\item Is the pool's design at fault or the attacker's conduct?
\item Should a pool accept thin \omterm{def:m3:blockchains-for-traders:key}{tokens} as collateral at all?
\item How should a supplier choose which pools to supply?
\item State the \emph{named result}: the pump that borrows the pool dry, and the attacker's cost against the gain.
\item In one sentence: what is an \omterm{def:m3:blockchains-for-traders:contract}{oracle} to a \omterm{def:m3:defi-credit-and-leverage:pool}{lending pool}?
\end{enumerate}
\end{problem}

\section{Interview questions}

\begin{interviewq}[$\star$ \iqroles{trader}]\label{iq:m3:defi-credit-and-leverage:1}
What is a \omterm{def:m3:defi-credit-and-leverage:ltv}{health factor}, and what happens when it falls below one?
\end{interviewq}

\begin{interviewq}[$\star$ \iqroles{developer}]\label{iq:m3:defi-credit-and-leverage:2}
How does a \omterm{def:m3:defi-credit-and-leverage:flash}{flash loan} work, and why can it only exist on a \omterm{def:m3:blockchains-for-traders:chain}{blockchain}?
\end{interviewq}

\begin{interviewq}[$\star\star$ \iqroles{researcher}]\label{iq:m3:defi-credit-and-leverage:3}
Why do \omterm{def:m3:defi-credit-and-leverage:pool}{lending pools} use a kinked rate curve rather than a linear one?
\end{interviewq}

\begin{interviewq}[$\star\star$ \iqroles{risk}]\label{iq:m3:defi-credit-and-leverage:4}
How would you set the loan-to-value of a new collateral \omterm{def:m3:blockchains-for-traders:key}{token}?
\end{interviewq}

\begin{interviewq}[$\star\star$ \iqroles{trader, researcher}]\label{iq:m3:defi-credit-and-leverage:5}
Why did an \omterm{def:m3:defi-credit-and-leverage:algo}{algorithmic stablecoin} paying 20\% collapse, and what would you have watched?
\end{interviewq}

\begin{interviewq}[$\star\star\star$ \iqroles{developer, risk}]\label{iq:m3:defi-credit-and-leverage:6}
Design a liquidation bot that competes profitably without taking bad inventory risk.
\end{interviewq}
